\documentclass[11pt]{article}

\usepackage[margin=1in]{geometry}
\usepackage{amsmath,amssymb}
\usepackage{booktabs}
\usepackage{graphicx}
\usepackage{hyperref}
\usepackage{microtype}
\usepackage{xcolor}

\title{From Pruned Weight Matrices to Sparse Graphs:\\
An Empirical Comparison of Neural-Network Storage Formats}
\author{Pranavi Immanni\\
CS 180 Individual Study\\
San Jos\'e State University}
\date{November 2026}

\newcommand{\todo}[1]{\textcolor{red}{[TODO: #1]}}

\begin{document}
\maketitle

\begin{abstract}
Pruning can reduce the number of active neural-network weights, but storing a
binary mask over a dense matrix may not provide the corresponding storage or
execution benefits.  This paper presents a reproducible comparison of dense,
dense-plus-mask, Coordinate (COO), Compressed Sparse Row (CSR), and weighted
adjacency-list representations for pruned neural-network weight matrices.
We study a fully connected Fashion-MNIST classifier with global magnitude,
per-neuron top-k, and random pruning at sparsities from 50\% to 99\%.
The study measures storage, inference latency, memory, prediction agreement,
and energy when a reliable energy measurement is available.  We also define a
lossless conversion from rectangular weight matrices to weighted bipartite
adjacency lists.  The current analytical storage results show that CSR is
substantially smaller than dense storage at high sparsity, while the runtime
and energy conclusions require representation-level benchmarks.  \todo{Add
final adjacency-list, latency, memory, and energy results before submission.}
\end{abstract}

\section{Introduction}

Modern neural networks frequently contain weight matrices in which many
parameters can be removed with limited effect on predictions.  A common
implementation represents this sparsity with a dense tensor and a binary mask.
Although this preserves the original matrix interface, it can retain much of
the dense storage and computation cost.  Sparse formats instead store only
nonzero values and the indices needed to locate them.

This work studies whether a pruned neural-network weight matrix can be treated
as a weighted graph and evaluated through sparse representations.  In
particular, a rectangular matrix is converted into a directed bipartite graph:
each output row is a row node, each input column is a column node, and a
nonzero weight becomes a weighted edge from the row node to the column node.
This construction permits a direct comparison between a graph-oriented
adjacency list and standard sparse matrix formats.

The central research question is:
\begin{quote}
When the same pruned weights are stored in different representations, which
format provides the best trade-off among storage, inference efficiency, memory,
and prediction preservation?
\end{quote}

The contributions of this study are:
\begin{enumerate}
    \item an implementation of signed magnitude, random, and per-neuron top-k
    pruning with deterministic masks;
    \item a lossless conversion of rectangular weight matrices into weighted
    bipartite adjacency lists;
    \item a controlled comparison of dense, dense-plus-mask, COO, CSR, and
    adjacency-list representations using identical retained weights;
    \item correctness and output-equivalence tests for the sparse forward
    operations; and
    \item a reproducible storage-cost model and benchmark protocol.
\end{enumerate}

\section{Background and Related Work}

\subsection{Sparse matrix representations}

COO stores a value together with its row and column indices for every stored
entry.  It is simple to construct, but repeats row information for every
nonzero.  CSR stores values and column indices for every nonzero together with
one row-pointer array.  The row pointers remove the repeated row indices and
make row-wise operations convenient.  These trade-offs are also reflected in
standard sparse-array libraries, where COO is commonly used for construction
and CSR is useful for arithmetic and vector products \cite{scipycoo}.

An adjacency list stores, for each source node, a list of destination nodes and
weights.  The project implementation uses Python dictionaries and lists for
clarity.  Its analytical byte count is COO-like, but actual Python-object
memory is larger and must be measured separately.

\subsection{Pruning and truly sparse computation}

Magnitude pruning ranks parameters by $|w|$, so negative weights are treated
symmetrically with positive weights while their original signs are retained.
Per-neuron top-k pruning ranks weights independently within each output row.
Random pruning provides a seeded control condition.

Recent work on truly sparse neural networks emphasizes that masked dense
implementations may preserve dense execution, while sparse kernels can avoid
materializing dense computations \cite{snack2026}.  Our study differs in scope:
it compares portable representations, including adjacency lists, COO, and CSR,
for pruned fully connected layers and reports the trade-offs of a direct
reference implementation.

\subsection{Neural-network compression}

Neural-network compression includes pruning, quantization, factorization, and
distillation.  Deep Compression combines pruning with trained quantization and
Huffman coding, demonstrating that storage reductions can be much larger than
the reduction in the number of nonzero connections alone \cite{deepcompression}.
Our study deliberately isolates the representation question: all formats use
the same full-precision retained weights, so a difference in storage or
latency can be attributed to the representation rather than to quantization or
entropy coding.

Magnitude pruning is a simple post-training baseline, but the resulting
unstructured sparsity does not automatically translate into a hardware
speedup.  The Lottery Ticket Hypothesis studies whether sparse subnetworks can
retain the trainability of dense networks \cite{lotteryticket}.  Dynamic sparse
training methods such as RigL change the sparse connectivity during training
and target sparse-to-sparse computation \cite{rigl}.  These works focus on how
to find or train sparse subnetworks.  Our focus is downstream: given a fixed
pruned matrix, how should it be stored and evaluated?

\subsection{Sparse execution and hardware effects}

The distinction between logical sparsity and realized efficiency is central to
this paper.  A mask can make entries logically inactive while leaving dense
matrix multiplication and dense memory allocation in the execution path.
Conversely, a sparse representation introduces index reads, irregular memory
accesses, conversion costs, and format-specific kernel overhead.  Prior work
has shown that structured sparsity can be easier for hardware to exploit than
unstructured sparsity; for example, structured dynamic sparse training reports
different CPU and GPU behavior for sparse linear layers \cite{srigl}.

SNACK is a particularly relevant recent system because it provides a custom
COO-based GPU kernel and explicitly targets truly sparse execution
\cite{snack2026}.  It is a systems contribution with optimized CUDA kernels,
training support, and large-scale GPU experiments.  We do not claim to match
that scope.  Instead, our contribution is a smaller and more transparent
comparison that includes a Python adjacency-list representation and uses the
same pruned weights across all formats.  This makes the study useful as a
baseline for understanding where a graph-oriented representation helps or
fails before investing in specialized kernels.

\subsection{Positioning and distinction}

The novelty of this paper is not a new pruning criterion and not a new COO or
CSR data structure.  The intended contribution is an empirical, reproducible
comparison of standard sparse formats and a weighted bipartite adjacency list
for neural-network weight matrices.  The comparison combines four properties
that are often evaluated separately: structural storage, actual inference
latency, memory behavior, and output equivalence.

\begin{table}[ht]
\centering
\caption{Positioning relative to representative prior work.}
\label{tab:positioning}
\small
\begin{tabular}{p{0.22\linewidth}p{0.18\linewidth}p{0.18\linewidth}p{0.30\linewidth}}
\toprule
Work & Main object & Primary contribution & Relation to this study \\
\midrule
Deep Compression & Neural-network parameters & Pruning, quantization, and coding & We hold value precision fixed and isolate storage formats. \\
Lottery Ticket Hypothesis & Sparse subnetworks & Trainability and initialization of sparse tickets & We start from fixed trained weights and study inference representation. \\
RigL / structured DST & Sparse training & Dynamic connectivity and structured sparsity & We do not propose a training algorithm; we compare fixed sparse layouts. \\
SNACK & Sparse GPU layers & Custom sparse kernels and end-to-end GPU efficiency & We provide a portable reference comparison including adjacency lists. \\
This work & Pruned MLP weight matrices & Storage, execution, and equivalence comparison & Same masks and weights across dense, mask, COO, CSR, and adjacency-list formats. \\
\bottomrule
\end{tabular}
\end{table}

The distinction should be stated honestly in the final paper: this is a
systems-oriented empirical study and reproducibility baseline, not a claim
that adjacency lists replace optimized sparse GPU kernels in every regime.

\section{Research Questions}

\begin{description}
    \item[RQ1:] How does storage scale with sparsity for dense, masked dense,
    COO, CSR, and adjacency-list representations?
    \item[RQ2:] Do sparse representations preserve the predictions and raw
    outputs of the corresponding pruned dense model?
    \item[RQ3:] Which representation provides the lowest inference latency at
    different sparsities and batch sizes?
    \item[RQ4:] Do analytical storage savings appear in actual memory usage?
    \item[RQ5:] What energy differences can be measured on the available
    hardware without relying on estimates?
\end{description}

\section{Methodology}

\subsection{Model and data}

The primary model is a fully connected Fashion-MNIST classifier with layer
dimensions $784\rightarrow300\rightarrow100\rightarrow10$, commonly called
LeNet-300-100.  ReLU activations occur between linear layers.  The current
baseline artifacts report a test accuracy of 0.8887 for one trained model.
\todo{Document all three model seeds, dataset preprocessing, and checkpoint
identifiers in the final version.}

\subsection{Pruning conditions}

We evaluate global magnitude, per-neuron top-k, and seeded random pruning at
50\%, 80\%, 90\%, 95\%, and 99\% requested sparsity.  All storage formats use
the same mask for a given pruning condition.  This isolates representation
effects from differences in which weights were retained.

Global magnitude pruning removes one globally ranked set of weights across all
layers.  Top-k pruning removes
$\lfloor s d_{\mathrm{in}}\rfloor$ entries from every output row of width
$d_{\mathrm{in}}$.  Random pruning uses a fixed seed and a global permutation.
Ties use deterministic layer and row-major ordering.

\subsection{Representations}

The evaluated representations are:
\begin{itemize}
    \item \textbf{pruned dense}: a dense tensor with removed entries set to
    zero;
    \item \textbf{dense plus mask}: the original dense tensor multiplied by a
    Boolean mask during the forward pass;
    \item \textbf{COO}: values, row indices, and column indices;
    \item \textbf{CSR}: values, column indices, and per-layer row pointers;
    \item \textbf{adjacency list}: a weighted row-to-column bipartite list.
\end{itemize}

For a weight matrix $W\in\mathbb{R}^{m\times n}$, the bipartite mapping uses
row nodes $0,\ldots,m-1$ and column nodes $m,\ldots,m+n-1$.  Each nonzero
$W_{ij}$ creates the edge $i\rightarrow m+j$ with weight $W_{ij}$.  Forward
propagation iterates the outgoing edges of each row node and accumulates the
corresponding input values.

\subsection{Metrics}

For every condition and batch size, we report:
\begin{enumerate}
    \item test accuracy;
    \item maximum absolute difference from the pruned dense output;
    \item prediction match rate;
    \item median and 95th-percentile inference latency;
    \item structural storage bytes;
    \item peak runtime memory using the platform-appropriate profiler; and
    \item energy only when directly measured by hardware instrumentation.
\end{enumerate}

Conversion time is reported separately from repeated inference time.  The
storage model uses 4-byte values, 8-byte indices, and one byte per mask entry
unless otherwise noted.

\section{Results}

\subsection{Analytical storage}

The model contains 266,200 weights, so dense fp32 storage is 1,064,800 bytes.
Dense-plus-mask storage is 1,331,000 bytes.  For CSR, the three layers add
$301+101+11=413$ row-pointer entries.

\begin{table}[ht]
\centering
\caption{Analytical storage for the LeNet-300-100 weights. Values use fp32
weights and 64-bit indices.}
\label{tab:storage}
\begin{tabular}{lrrrrr}
\toprule
Sparsity & Dense & Dense+mask & COO & CSR & CSR / dense \\
\midrule
50\% & 1,064,800 & 1,331,000 & 2,662,000 & 1,600,504 & 150.3\% \\
80\% & 1,064,800 & 1,331,000 & 1,064,800 & 642,184 & 60.3\% \\
90\% & 1,064,800 & 1,331,000 & 532,400 & 322,744 & 30.3\% \\
95\% & 1,064,800 & 1,331,000 & 266,200 & 163,024 & 15.3\% \\
99\% & 1,064,800 & 1,331,000 & 53,240 & 35,248 & 3.3\% \\
\bottomrule
\end{tabular}
\end{table}

At 90\% sparsity, changing CSR indices from 64-bit to 32-bit reduces the
analytical cost to 214,612 bytes, or 20.15\% of dense fp32 storage.  Using
32-bit indices and fp16 values gives 161,372 bytes, or 15.16\%.

These values are structural accounting, not measurements of the full Python
object graph or process peak memory.  The generated comparison file is
`reports/cost_model_vs_measured.csv`.

\subsection{Prediction preservation}

The conversion and list-matvec tests pass for rectangular matrices, negative
weights, empty rows, symmetric bipartite graphs, and three-layer MLP forward
propagation.  The format-level Fashion-MNIST results are included in the
project's benchmark artifacts.  \todo{Add the final table of test accuracy,
prediction agreement, and maximum output error.}

\subsection{Latency, memory, and energy}

\todo{Add benchmark tables and plots after the adjacency-list experiment is
run using the same trained checkpoint and test set as the existing formats.}

\section{Discussion}

The storage model predicts that CSR becomes smaller than dense storage once
sparsity is sufficiently high, while COO pays for two indices per nonzero.
Dense-plus-mask remains larger than dense under the byte accounting used here
because every weight and every mask entry is retained.  These structural
predictions do not imply a runtime speedup: sparse kernels can have indexing,
launch, and Python-object overheads, especially for small matrices or small
batches.  The benchmark therefore treats storage and execution as separate
questions.

\todo{Discuss the measured crossover points, batch-size effects, and the
difference between Python adjacency lists and compiled sparse kernels.}

\section{Limitations}

This study uses one fully connected architecture and one dataset.  The
adjacency-list reference implementation is written in Python and should not be
interpreted as a production GPU sparse kernel.  Actual energy results depend on
available hardware instrumentation.  Finally, analytical storage estimates
exclude allocator overhead, object headers, and temporary tensors; those costs
must be reported separately from structural bytes.

\section{Conclusion}

This paper establishes a controlled framework for converting pruned neural
network weight matrices into sparse graph representations and comparing them
with standard sparse formats.  The current storage analysis shows substantial
CSR savings at high sparsity while preserving the dense model's parameter
values in the sparse representations.  \todo{Replace this preliminary
conclusion with the final latency, memory, energy, and adjacency-list findings.}

\bibliographystyle{plain}
\bibliography{references}

\end{document}
